\section{Limitations}
\label{sec:limitations}

Our interventions are controlled camera changes, not a prevalence study of
deployed preprocessing bugs.  The GT evidence covers 13 ETH3D scenes, 22 DTU
scans, three or four views, and two model families.  It does not establish the
same severity for other view counts, resolutions, training preprocessors,
domains, or reconstruction systems.  The support-preserving control makes the
coordinate effect causal under its fixed canvases, but it was designed after
the ETH3D crop study; its chronology is disclosed and it should be replicated
independently.  Connecting the audit to logged failures in a production or
user-captured pipeline remains important future evidence.

Cropping changes coordinates and visible support together.  Canonical repair
holds crop-supported RGB samples fixed while changing canvas placement, but it
adds a fixed fill and another interpolation step and cannot reconstruct removed
pixels.  Consistent with this boundary, it repairs camera orientation but
worsens ETH3D depth for both models.  DTU offers no matched depth maps under our
protocol, and strict observation/20-mm point filtering leaves some DUSt3R point
directions undefined; we refuse a subset bootstrap.  The DTU point-map metric
uses camera-pose alignment and a deterministic cap and is not the official MVS
leaderboard protocol.

Cross-transform disagreement passes the held-out AUROC gate for both models,
but it only ranks the frozen intervention cases and is not a calibrated
probability.  It strongly beats VGGT's native score yet ties DUSt3R native AUROC,
so universal superiority is unsupported.  Likewise, three-fill consensus does
not beat the analytic baseline across both models.  Finally, DUSt3R carries a
non-commercial code and checkpoint license; we distribute only our adapter,
protocol, hashes, and lightweight evaluation evidence.
